We compare CMDP rollout, which we compute via quasi-Monte Carlo integration, to EI and EIpu.
We use a GP with the Mat\'ern-$5/2$ ARD kernel to model both the objective and the cost function\footnote{We use a log-warped GP to model positive cost.}, and learn hyperparameters via maximum likelihood estimation.
When rolling out acquisition functions, we use L-BFGS-B using $5$ restarts, selected by evaluating the acquisition on a Latin hypercube of $10d$ points and picking the five best as starting points.
When comparing different replications we first need to interpolate the objective function values onto a set of discrete costs.
Given these interpolated value, we plot the mean with one standard deviation. Code to reproduce our experiments is found at \url{https://github.com/ericlee0803/lookahead_release}.

\subsection{Synthetic Problem}
\begin{figure}[!ht]
    \centering
    \includegraphics[width=0.48\textwidth]{figures/costrollout_synthetic2.pdf}
    \caption{In this example, we examine a carefully chosen example showcasing the strength of the rollout approach. We consider the a multimodal objective whose most expensive point is the global minimum. EIpu performs worse than EI, and both tend to get stuck in cheaper, local minimum. Our rollout policy for horizons $2$ and $4$ performs better than both EI and EIpu}
    \label{fig:costrollout2d}
\end{figure}

In Figure~\ref{fig:costrollout2d}, we examine a carefully chosen synthetic example showcasing the strength of the rollout approach. We consider the cost-constrained optimization problem
\begin{align*}
    f(\bx) &= 10 \|\bx\|_2 \sin(2\pi \|\bx\|_2), \\
    c(\bx) &= 10 - 5\|\bx\|_2,
\end{align*}
in the domain $[-1, 1]^2$, and a budget of $150$. The cost function has been designed so that its maximum aligns with the minimum of the objective.
As we motivated earlier, EIpu struggles with these types of problems.
We run BO with EI, EIpu, and rollout with our base policy $50$ times and plot the results.
This is seen on the right, in which EIpu (green) performs worse than EI (blue).
However, rollout of our base policy, for horizons two and four in pink and red respectively, performs much better than both.

\subsection{Hyperparameter Optimization}
\begin{figure*}[!ht]
    \centering
    \includegraphics[width=\textwidth]{figures/HPO_results.pdf}
    \caption{We compare the classification error among EI, EIpu, and our cost-constrained rollout for horizons $2$ and $4$. Rollout performs better than both EI and EIpu. Shaded areas represent one standard deviation around the mean. }
    \label{fig:hpo_results}
\end{figure*}
We compare rollout performance to EI and EIpu on HPO of three different models: $k$-nearest neighbors (KNN), decision trees, and random forests, with budgets of $800$, $15$, and $15$ seconds respectively. These are relatively small problems, chosen due to the number of replications required to show statistical significance.
All models are trained for $50$ replications on the OpenML w2a dataset~\citep{OpenML2013}. We compare the competing algorithms in terms of the best classification error achieved on the validation set.

\textbf{$k$-nearest neighbors:} The $k$-nearest neighbors algorithm is a class of methods used for classification of either spatially-orientated data or data with a known distance metric (i.e., data embedded in a Hilbert space).
We consider a $5d$ search space: dimensionality reduction percentage in $[1\mathrm{e}{-6}, 1.0]$ (log-scaled), type in \{Gaussian, Random\}, neighbor count in $\{1, 2, \ldots, 256\}$, weight function in \{Uniform, Distance\}, and distance in \{Minkowski, Cityblock, Cosine, Euclidean, L$1$, L$2$, Manhattan\}. We one-hot encode categorical variables.\looseness-1

\textbf{Decision Trees:} Decision trees are popular predictive models used in statistics, data mining, and machine learning.
In the case of classification, leaves represent class labels and paths represent sets of features that lead to those class labels.
During training, a tree is built by splitting the source set into subsets which constitute the successor children.
The splitting is based on a threshold that maximizes some notion of information gain such as entropy.
The depth of a decision tree is pre-specified.
We consider a $3d$ search space: tree depth in $\{1, 2, \ldots, 64\}$, tree split threshold in $[0.1, 1.0]$ log-scaled, and split feature size in $[1\mathrm{e}{-3}, 0.5]$ (log-scaled).

\textbf{Random Forests:} A random forest is a set $k$ of decision trees, and classifies based off the plurality decision generated from all its trees ---this technique is known as \textit{bagging}, and improves robustness in the classification algorithm.
We consider a $3d$ search space: number of trees in $\{1, 2, \ldots, 256\}$, tree depth in $\{1, 2, \ldots, 64\}$, and tree split threshold in $[0.1, 1.0]$ (log-scaled).

\subsection{Sensor Set Selection}

The sensor set selection problem~\citep{garnett2010bayesian} seeks to improve the predictive accuracy, as measured by the root mean squared error (RMSE), of a physical sensor network.
We denote a sensor network's configuration of $m$ sensors in $d$-dimensional space as $\mathbf{X}\in\R^{m \times d}$.

This configuration must be manually adjusted each time it is updated.
This is typically assumed to have uniform cost; we modify the problem to consider a \textit{sensor adjustment cost}.
We assume this cost is correlated with the distance each sensor in the network has to move; for simplicity, we assume the cost of any new sensor configuration $\mathbf{X}'$ is proportional to the straight-line distance $d(\mathbf{X}, \mathbf{X}') = \|\mathbf{X} - \mathbf{X}'\|_F$.
This is an example of a CMDP whose cost is not state-independent; our cost now depends on the prior state.

We consider a small sensor set selection problem using ten sensors.
Our objective is RMSE of the sensor predictions against ground truth weather data taken from UK Meteorological Office MIDAS surface stations~\citep{midas2012}.
The time budget is twenty years.

\subsection{Analysis}
We plot the resuts of our HPO and sensor set selection experiments in Figure~\ref{fig:hpo_results}, and find that CMDP rollout generally outperforms both EI and EIpu.
In this section we discuss key insights gained over the course of experimentation.

\textbf{Cost Modeling:} We found the cost function to be simpler to model than the objective function.
In practice, the cost may only depend on a few key parameters (e.g., tree depth).
Thus, using a vanilla GP to model the cost is inefficient ---a tailored (parametric) cost model or a GP that incorporates parameter importance into its lengthscale priors will likely lead to better results~\citep{lee2020cost, guinet2020}.

\textbf{Search Space Sensitivity:}
EIpu's performance depends on the correlation between objective value and cost.
Unsurprisingly, this correlation often depends on the search space in practice.
For example, assume a decision tree of depth $d$ achieves maximal classification error and that its training cost increases with depth.
(\textbf{i}) If the search space is $[1, d\,]$, the maximum will be the most expensive point and EIpu will perform poorly; (\textbf{ii}) if the search space is $[1, 2d\,]$, the maximum will be have middling cost and EIpu will perform moderately well; (\textbf{iii}) if the search space is $[1, 10d\,]$, the maximum will have cheap cost and EIpu will perform very well.
In our experiments, we found CMDP rollout to be more robust to the shape of the cost surface.
This is expected, as a CMDP optimal policy selects the point that maximally reduces the objective function given the cost constraint.
